\chapter{Jazz: Writing a Conversation into Harmony}
\label{ch:genre-jazz-deep}

\classquote{大乐必易，大礼必简。}{《礼记·乐记》}{Great music is grounded in simplicity; great ritual rests on restraint. Beneath jazz's surface complexity sits a minimal set of conversational rules.}

\section{What jazz is: listen to how they play ``together''}
\label{sec:jazz-what}

\term{jazz}{jazz} is not a scale, and it is not ``complicated chords''. It is closer to an agreement about how to play together: one musician states a theme, another answers; the rhythm section walks steadily underneath while the front line builds sentences in the moment. The most moving moments on a jazz record are often not written notes at all, but a shared decision made on the same beat.

To hear jazz, start with four handles:

\begin{longtable}{p{0.16\textwidth}p{0.32\textwidth}p{0.38\textwidth}}
\toprule
Handle & What it feels like & The writing question \\
\midrule
Swing & Eighths lean ``long--short''; the body wants to move & Is this swing really pushing the offbeat back, or is it even? \\
Walking bass & One note per beat, laying steps on the floor & Does the bass only play roots, or connect them with passing tones? \\
Call and response & One line asks, another answers, like people talking & Do comping and solo leave room for an answer? \\
Improvisation & New sentences unfold after the theme, different each time & Does the solo develop a motif over the chords, or just run a scale? \\
\bottomrule
\end{longtable}

The opposite of jazz is not ``non-standard''; it is ``no conversation''. A solo that is only a chord progression with no response is like talking at a wall. Jazz professionalism comes from listening: the rhythm section answers a soloist's push, and the soloist fills the space that was left open.

\section{From New Orleans to the world: how jazz got here}
\label{sec:jazz-history}

Jazz absorbed blues, ragtime, marching bands, gospel and Latin rhythms, then reorganised them into a language built for improvisation. The table follows one thread: when the setting changed how people played, the instrumentation changed with it.

\begin{longtable}{p{0.16\textwidth}p{0.34\textwidth}p{0.42\textwidth}}
\toprule
Period & Key shift & Effect on arrangement \\
\midrule
\endfirsthead
\toprule
Period & Key shift & Effect on arrangement \\
\midrule
\endhead
1890s--1910s & New Orleans brass bands, funerals and parades; blues and ragtime & Front-line horns improvise over a rhythm section; collective improvisation dominates \\
1920s & The Jazz Age, records and dance halls; Chicago and New York & The small-band lineup takes shape: trumpet, clarinet/sax, piano, bass, drums \\
1930s--1940s & Big bands and swing & Sections of trumpets, trombones and saxophones plus rhythm; the arranger becomes a central role \\
1940s & Bebop & Faster tempos, denser chords; back to small bands, spotlighting the soloist \\
1950s & Cool jazz, hard bop and modal jazz & Timbres narrow or get bluesier; \term{modal}{modal} thinking reduces chord density and gives soloists room \\
1960s & Free jazz and the roots of fusion & Looser structure, electric instruments enter; metre and tonality are reopened \\
1970s--1990s & Fusion, electric sounds and crossover & Electric bass and synthesizers enter; jazz harmony mixes with pop arrangements \\
2000s--now & Global touring, streaming and cross-genre work & Young musicians know jazz, hip-hop, electronic and classical at once; genre and lineup boundaries soften \\
\bottomrule
\end{longtable}

Jazz is also a travelling language: wherever it went, it merged with local music---Brazilian bossa nova, Cuban Afro-Cuban, European improvised theatre, Japan's jazz-cafe culture. For many Chinese listeners, jazz first arrives as ``a piano trio in a cafe'', and only later reveals its blues roots and improvisational core.

\section{Instrumentation: the rhythm section and the front line}
\label{sec:jazz-craft}

Jazz instrumentation fits in one sentence: \textbf{the rhythm section owns time and harmony; the front line owns melody and story}. Write that into your track names and it beats stacking ten presets.

\subsection{Typical lineup and two roles}
\label{sec:jazz-instrumentation}

\begin{longtable}{p{0.16\textwidth}p{0.25\textwidth}p{0.25\textwidth}p{0.20\textwidth}}
\toprule
Layer & Common instruments/tracks & Main job & Common mistake \\
\midrule
\endfirsthead
\toprule
Layer & Common instruments/tracks & Main job & Common mistake \\
\midrule
\endhead
Rhythm section & Piano/guitar, bass, drums & Set time, carry harmony, leave space & All three fight for the same register \\
Front line & Trumpet, saxophone, trombone, voice & Give the theme, solo, tell the story & The solo runs from start to finish with no space \\
Colour & Vibraphone, strings, organ & Thicken texture, change section colour & Colour layer covers the walking bass' low end \\
\bottomrule
\end{longtable}

The simplest starting point is four parts: piano, bass, drums, plus one melodic front-line voice (trumpet or saxophone). Get the ``rhythm section + front line'' relationship working first, then decide whether to add a second horn or a vibraphone. Track count is not jazz depth; conversation is.

\subsection{Groove: swing is not ``write it long--short''}
\label{sec:jazz-groove}

\term{Swing}{swing} pushes the second half of each eighth pair later, creating a ``long--short'' feel. The key insight: swing is a \textbf{performance feel}, not a mechanical rewrite of note values. In Yeban you can use swing quantize or nudge the offbeats by hand, but the feel comes mostly from the ride cymbal and the drummer's accent pattern:

\begin{itemize}
\item \textbf{Ride cymbal}: an ongoing ``ding--ding--da'' swing pattern; the clock of jazz time.
\item \textbf{Snare}: soft taps on 2 and 4, not rock-level hits.
\item \textbf{Kick drum}: sparse feathering, staying out of the melody's way.
\item \textbf{Comping}: piano/guitar voice chords on key beats and leave the rest empty---the answer in a call-and-response.
\end{itemize}

If every beat is equally full, swing loses its elasticity. A rhythm section's craft is not ``hit every beat'' but ``stay steady when walking, and truly leave space when space is called for''.

\subsection{Harmony: ii--V--I and the two kinds of space}
\label{sec:jazz-harmony}

The core grammar of jazz harmony is \term{ii--V--I}{ii-V-I}: a minor ii, a dominant V and a major I resolving. Blues uses I--IV--V and allows \term{blue notes}{blue-note} (bent third, fifth and seventh). The two grammars offer different writing spaces:

\begin{itemize}
\item \textbf{Functional harmony (bebop etc.)}: dense chords; the soloist chases changes---good for practising speed and control.
\item \term{Modal}{modal} jazz: few chords, long stays on one mode; the soloist works within one area---good for developing motifs.
\end{itemize}

A practical route in Yeban: practise modal jazz first for ``few chords, long lines'', then ii--V--I for ``fast chords, short motifs''. The former teaches development, the latter teaches keeping up.

\subsection{Melody: short theme, long solo}
\label{sec:jazz-melody}

A jazz \term{head}{head} is usually a memorable 8, 16 or 32-bar phrase; the real expansion is the \term{improvisation}{improvisation} after it. Two rules for writing a head:

\begin{enumerate}
\item Short enough to remember and vary---not a long pop-style line.
\item Leave one developable motif (a rhythm or two or three notes) the solo can build on.
\end{enumerate}

A solo is not a scale run top to bottom. A good solo repeats a motif, moves it to new chords, changes its rhythm and returns to the theme at the end---so the listener hears the old theme inside the new sentences.

\subsection{Sound design: give the solo somewhere to stand}
\label{sec:jazz-sound}

Much of jazz ``arranging'' is \textbf{dynamics}: who is loud, who steps back. Three actionable rules:

\begin{itemize}
\item \textbf{Rhythm section drops density when a solo enters}: piano voices fewer chords, drummer trades hits for ride and soft snare, bass keeps walking without filling.
\item \textbf{Call and response needs real space}: after comping a chord, give the soloist half a beat to a beat to answer.
\item \textbf{One owner of the low end}: walking bass owns the lows; don't add a competing bass when piano left hand doubles it, and keep vibraphone low bars out of the way.
\end{itemize}

\section{Four pieces, arranged}
\label{sec:jazz-cases}

The four pieces below are arrangement analyses only---no lyrics, no full melodies. Work information is sourced from album releases and official pages; links are in Further Reading.

\subsection{``So What'': less-is-more modal model}
\label{sec:jazz-case-so-what}

Miles Davis' ``So What'' appears on the 1959 album \emph{Kind of Blue}. A modal-jazz landmark, it stays in two modes almost throughout and leaves the soloists enormous space.

\begin{itemize}
\item \textbf{Harmony}: not dense ii--V--I, but long stays on one mode; soloists develop motifs ``horizontally'' instead of chasing chords ``vertically''.
\item \textbf{Theme}: a minimal descending motif, stated by trumpet and answered by saxophone---the theme itself is a call and response.
\item \textbf{Rhythm section}: sparse piano comping, steady walking bass, ride keeping time; the whole group ``steps back'' cleanly.
\item \textbf{Arrangement logic}: subtraction creates space so every soloist's line is audible.
\end{itemize}

Takeaway: if your solos always sound like scale runs, cut to two chords, stretch the time, and force yourself to develop a motif instead of switching chords.

\subsection{``Take the 'A' Train'': big-band section writing}
\label{sec:jazz-case-a-train}

Duke Ellington and Billy Strayhorn's ``Take the 'A' Train'' is a 1940s big-band swing signature. It shows how arranging turns a group of instruments into one layered whole.

\begin{itemize}
\item \textbf{Lineup}: trumpets, trombones and saxophones in sections, answering the rhythm section; sections take turns carrying theme and response.
\item \textbf{Timbre}: mutes and registers distinguish sections, so ``who is speaking'' is instantly clear.
\item \textbf{Groove}: swing drive comes from the whole band's aligned accents, not one player's fill.
\item \textbf{Arrangement logic}: write section functions first, melody second; every group has a job.
\end{itemize}

Takeaway: even with only five MIDI tracks, mimic big-band ``section feel''---group instruments and let groups answer each other instead of playing one melody together.

\subsection{``My Favorite Things'': turning a theme into solo building blocks}
\label{sec:jazz-case-favorite-things}

John Coltrane's 1960 recording of ``My Favorite Things'' shows how a simple theme becomes the material for a long solo.

\begin{itemize}
\item \textbf{Theme}: a memorable, regular melody---both the theme and the solo's ``building blocks''.
\item \textbf{Solo}: Coltrane breaks the motif apart, transposes it, thickens it and reassembles it; the listener keeps hearing the theme's shadow.
\item \textbf{Rhythm section}: piano and bass hold a steady cycle through the long solo, giving a reliable floor.
\item \textbf{Arrangement logic}: not ``theme + free improv'', but ``theme → transformation of theme → return of theme''.
\end{itemize}

Takeaway: before soloing, write down the two or three notes you most want to develop; bring every solo back to that motif so listeners can follow.

\subsection{``Change of the Guard'': jazz in a contemporary lineup}
\label{sec:jazz-case-change-of-guard}

Kamasi Washington's ``Change of the Guard'', from the 2015 album \emph{The Epic}, puts hard-bop language, funk grooves and contemporary production in one lineup---a door for younger listeners into jazz.

\begin{itemize}
\item \textbf{Lineup}: saxophones, trumpets and trombones up front; a rhythm section of electric keys/bass and drums with hip-hop/funk feel.
\item \textbf{Layers}: big-band section thinking combined with modern overdubs---thick but clear.
\item \textbf{Groove}: both swinging and funk-accented, showing jazz feel can coexist with contemporary rhythm.
\item \textbf{Arrangement logic}: keep jazz's conversation and solo core, swap the outer layer for timbres and energy younger ears know.
\end{itemize}

Takeaway: jazz need not be ``old-school''. Keep the conversation and improvisation, change to a timbre and groove you know, and you'll write a more personal voice.

\section{Yeban sketch: an eight-bar jazz trio}
\label{sec:jazz-yeban}

An eight-bar sketch you can rebuild in Yeban, using ``piano trio + one front-line melody'' to run a jazz conversation.

\begin{longtable}{p{0.10\textwidth}p{0.22\textwidth}p{0.58\textwidth}}
\toprule
Bars & Section & Arrangement moves \\
\midrule
\endfirsthead
\toprule
Bars & Section & Arrangement moves \\
\midrule
\endhead
1--4 & Head (theme) & Piano: right hand writes an 8-bar theme motif, left hand voices chords only on 1/3; Bass: one-note-per-beat walking line (root--passing--fifth--passing); Drums: swing eighths on ride, soft snare on 2/4; trumpet/sax: call-and-response with the piano theme, leaving two beats at the phrase end \\
5--8 & Solo & Front-line develops bars 1--4's motif over ii--V--I; piano comping drops density, chords only on key beats; bass keeps walking without fills; drums trade snare for ride and soft hi-hat; return to the theme motif in the last two bars \\
\bottomrule
\end{longtable}

Build order:

\begin{sloppypar}
\begin{enumerate}
\item Create a medium-tempo (about 120 BPM), 4/4 swing project and build tracks as above: Melody (MIDI), Piano (MIDI), Bass (MIDI), Drums (MIDI), Master bus.
\item Write the bass' eight-bar walking line first: one note per beat, passing tones connecting chord roots; solo-listen to confirm it ``walks'' rather than just marking roots.
\item Write drums: swing eighths on ride, soft snare on 2/4, sparse kick; enable swing quantize so eighths really lean.
\item Write the theme: target notes on ii--V--I chord tones, two beats left at the phrase end; give the theme a developable two-note motif.
\item Write piano comping: chords only on key beats, the rest empty---answering the front line.
\item In the solo section thin the piano and drum density; open the piano roll (F7) and move the motif onto new chords, repeating and developing it.
\item Open the device stack (F3) to confirm the front line and piano aren't fighting for the same register; check levels (F8) so the bass alone owns the low end.
\item Return to the theme motif in the last two bars, save and export an eight-bar listen; if you can hear ``ask--answer--develop--return'', it works.
\end{enumerate}
\end{sloppypar}

\section{Exercises: three routes, different trade-offs}
\label{sec:jazz-practice}

\begin{exercise}
\textbf{One theme, three solo openings.} Write an 8-bar theme, then three solos: A develops a single motif rhythmically; B moves the motif to each chord; C runs chord tones before returning to the theme. Compare coherence, surprise and comprehensibility. A is most coherent but risks monotony; B feels most structured; C is easiest to follow but risks sounding like an etude.
\end{exercise}

\begin{exercise}
\textbf{Write the swing, then undo it.} In Yeban, write a passage of even eighths and play it; then enable swing quantize (or nudge offbeats by hand) and play again. Record which one makes you want to move. Finally offset the bass' and ride's swing slightly and hear where the elasticity comes from.
\end{exercise}

\begin{exercise}
\textbf{Four-track big band.} Use only four MIDI tracks to fake big-band section feel: split them into an ``asking group'' and an ``answering group'' and have them alternate the same motif in different registers. Compare unison with grouped call-and-response; use no extra tracks---only division of labour.
\end{exercise}

\section{Further reading}
\label{sec:jazz-reading}

\begin{sloppypar}
\begin{itemize}
\item \textbf{Wikipedia, ``Jazz'':} \url{https://en.wikipedia.org/wiki/Jazz}. A genre map of jazz's branches---start here for the big picture.
\item \textbf{Wikipedia, ``Kind of Blue'':} \url{https://en.wikipedia.org/wiki/Kind_of_Blue}. Miles Davis' 1959 album entry, with the background and modal method of ``So What''; for audio, search the title on the artist's official channel.
\item \textbf{Wikipedia, ``Take the 'A' Train'':} \url{https://en.wikipedia.org/wiki/Take_the_%22A%22_Train}. The Duke Ellington/Billy Strayhorn composition entry---lineup and era.
\item \textbf{Wikipedia, ``My Favorite Things (John Coltrane album)'':} \url{https://en.wikipedia.org/wiki/My_Favorite_Things_(John_Coltrane_album)}. Coltrane's 1961 album entry, explaining how the title track turns a Broadway tune into a modal long solo.
\item \textbf{Wikipedia, ``The Epic (Kamasi Washington album)'':} \url{https://en.wikipedia.org/wiki/The_Epic_(album)}. The 2015 album entry, with lineup and release details for ``Change of the Guard''.
\end{itemize}
\end{sloppypar}

The links above were checked against entry names and release data on 27 November 2026 (album/composer/date are in each entry's text); Wikipedia URLs are stable---if your region blocks them, search the same title in a local language edition. For timbre and improvisation, listen on the artists' official channels by title.

\section{Takeaways}
\label{sec:jazz-takeaways}

\begin{itemize}
\item Jazz's core is ``conversation'': theme states, others answer, rhythm section supports, solo develops.
\item Swing is a performance feel, not long--short note values; ride cymbal and time feel matter more than scales.
\item Remember ``rhythm section + front line'': rhythm owns time and harmony, front line owns melody and story.
\item ii--V--I and modal harmony give two spaces: one for keeping up with change, one for developing motifs.
\item Keep the head short and memorable with a developable motif; let the solo return to the theme, not run scales.
\item Drop rhythm-section density when a solo enters, leave real space for call and response, and let walking bass alone own the low end.
\item Build the Yeban sketch in the order ``bass--drums--theme--comping--solo--register/level check--return''.
\end{itemize}

\begin{helpmapping}
\begin{longtable}{p{0.45\textwidth}p{0.45\textwidth}}
\toprule
Section & Help anchor \\
\midrule
What jazz is & \texttt{arrangement-genre-jazz-what} \\
From New Orleans to the world & \texttt{arrangement-genre-jazz-history} \\
Instrumentation: rhythm section and front line & \texttt{arrangement-genre-jazz-craft} \\
Four pieces, arranged & \texttt{arrangement-genre-jazz-cases} \\
Yeban sketch & \texttt{arrangement-genre-jazz-yeban} \\
Exercises & \texttt{arrangement-genre-jazz-practice} \\
Further reading & \texttt{arrangement-genre-jazz-reading} \\
\bottomrule
\end{longtable}
\end{helpmapping}
